% \documentclass[polish, 8pt]{article}
\documentclass[polish, 10pt]{extarticle}
\usepackage[top=2cm, bottom=2cm, left=3cm, right=3cm]{geometry}
\usepackage{graphicx} % Required for inserting images

\usepackage[polish]{babel}
\usepackage{polski}
\usepackage{placeins}
\usepackage{dirtytalk}
% Set page size and margins
% Replace `letterpaper' with `a4paper' for UK/EU standard size
\usepackage{csquotes}
% Useful packages
\usepackage{float}
\usepackage{amsmath}
\usepackage{graphicx}
\usepackage{hyperref}
\usepackage{graphicx}
\usepackage{listings}
\usepackage{xcolor}
\usepackage{amssymb}
\usepackage{enumitem}
\usepackage{float}
\definecolor{codegreen}{rgb}{0,0.6,0}
\definecolor{codegray}{rgb}{0.5,0.5,0.5}
\definecolor{codepurple}{rgb}{0.58,0,0.82}
\definecolor{backcolour}{rgb}{0.95,0.95,0.92}

\lstdefinestyle{mystyle}{
    backgroundcolor=\color{backcolour},   
    commentstyle=\color{codegreen},
    keywordstyle=\color{magenta},
    numberstyle=\tiny\color{codegray},
    stringstyle=\color{codepurple},
    basicstyle=\ttfamily\footnotesize,
    breakatwhitespace=false,         
    breaklines=true,                 
    captionpos=b,                    
    keepspaces=true,                 
    numbers=left,                    
    numbersep=5pt,                  
    showspaces=false,                
    showstringspaces=false,
    showtabs=false,                  
    tabsize=2
}
\lstset{style=mystyle}
\hypersetup{
    colorlinks=true,
    linktoc=all     %set to all if you want both sections and subsections linked
}
\usepackage{mdframed}
\usepackage{titling}
\setlength{\droptitle}{-2cm} 

\begin{document}

\title{Rachunek Prawdopodobieństwa i Statystyka -- Kolokwium I}
\author{}
\date{}

\maketitle
\vspace{-2cm}

\section*{Zadanie 1 (20 pkt)}
Sędziowie prowadzący popularny program telewizyjny \textit{Taniec ze Statystykami}, wykorzystują swoją władzę i przyjmują łapówki od uczestników. W każdym odcinku dany uczestnik może albo zostać w programie, albo zostać wyrzucony.
\begin{itemize}
\item Jeśli uczestnik daje łapówki, na pewno zostanie w programie.
\item Jeśli uczestnik nie daje łapówek, zostanie wyrzucony z programu w $\frac{2}{3}$ przypadków.
\end{itemize}
Załóżmy, że 25\% uczestników oszukuje. Ci sami uczestnicy dają łapówki w obu rundach, tj. jeśli uczestnik daje łapówki w pierwszej rundzie, robi to również w drugiej (i vice versa).
\begin{itemize}
\item[a)] Jeśli wybierzesz losowego uczestnika, który został w programie podczas pierwszego odcinka, jakie jest prawdopodobieństwo, że dawał łapówki sędziom?
\item[b)] Jeśli wybierzesz losowego uczestnika, jakie jest prawdopodobieństwo, że zostanie w programie podczas obu pierwszych odcinków?
\item[c)] Jeśli wybierzesz losowego uczestnika, który został w programie podczas pierwszego odcinka, jakie jest prawdopodobieństwo, że zostanie wyrzucony w drugim odcinku?
\end{itemize}

\section*{Zadanie 2 (15 pkt)}
\begin{itemize} 
\item Podaj definicję miary probabilistycznej oraz przynajmniej 4 z (podanych na wykładzie) jej podstawowych własności.
\item Udowodnij, że dla dowolnych dwóch zdarzeń $A, B \in \Sigma$ spełniona jest nierówność:
$$ P(A, B) \geq P(A) + P(B) - 1 $$
\end{itemize}

\section*{Zadanie 3 (30 pkt)}
Funkcja $f$ jest funkcją gęstości prawdopodobieństwa zmiennej losowej $X$:
\[
f(x) = 
\begin{cases} 
ax + bx^2 & \text{dla } 0 \leq x \leq 1, \\
0 & \text{w przeciwnym przypadku}.
\end{cases}
\]
Wiedząc, że wartość oczekiwaną $\mathbb{E}(X) = \frac{7}{12}$, wyznacz współczynniki $a$ i $b$, oblicz wariancję i modę, wyznacz dystrybuantę $F_X(x)$. Oblicz $P(|X| < \frac{1}{2})$.


\section*{Zadanie 4 (20 pkt)}
Na dwóch lotniskach w ciągu 30 minut mają wylądować niezależnie
dwie pary samolotów. W każdej parze jest maszyna typu A i maszyna typu B.
Samolot typu A zajmuje w czasie lądowania pas na 7 minut, a samolot typu B
na 9 minut. Jakie jest prawdopodobieństwo tego, że na przynajmniej jednym
z lotnisk samoloty nie będą na siebie czekać?

\section*{Zadanie 5 (15 pkt)}
Niech $X \sim U(a, b)$ dla $0 < a < b$. Wyznacz rozkład zmiennej losowej $Y = \ln X$. Dla jakich wartości $c, d \in \mathbb{R}$ spełniona jest poniższa równość?
$$P(c < Y \leq d) = f_Y(d) - f_Y(c)$$
% \section*{Zadanie 5 (20 pkt)}
% $X_1$, $X_2$ i $X_3$ to zmienne o rozkładzie jednostajnym: $X_1, X_2, X_3 \sim U(0, 1)$. Wyznacz rozkład:
% \begin{itemize}
%     \item zmiennej $Y = \sqrt{X1}$,
%     \item zmiennej $Z = \max(X_2, X_3)$.
% \end{itemize}
% Jaki jest związek między tymi zmiennymi? Zinterpretuj wynik.

\end{document}

% \section*{Zadanie 3}
% Rozważ następującą grę. Gracz rzuca kostką dokładnie jeden raz. Wygrana wynosi 1\$ za każdą \textit{kropkę} na wyrzuconej ściance kostki. Zakładając, że kostka jest uczciwa, na jakim poziomie należy ustalić cenę biletu za udział w tej grze?
